\section{问题表述与符号}\label{sec:formulation}

令 $q=[q_0,q_1,q_2,q_3]^\top\in\mathbb{R}^4$ 为单位四元数，$R(q)\in SO(3)$ 为相应的旋转。我们使用齐次提升
\begin{equation}
x_R(q)=\big[\|q\|^2,\ \mathrm{vec}(R(q))^\top\big]^\top\in\mathbb{R}^{10},
\label{eq:lift}
\end{equation}
以及 Veronese 向量
\begin{equation}
\begin{split}
v_2(q)=\big[&q_0^2,\ q_1^2,\ q_2^2,\ q_3^2,\ \sqrt2q_0q_1,\ \sqrt2q_0q_2,\ \sqrt2q_0q_3,\\
&\sqrt2q_1q_2,\ \sqrt2q_1q_3,\ \sqrt2q_2q_3\big]^\top\in\mathbb{R}^{10}.
\end{split}
\label{eq:veronese}
\end{equation}
QPEP 代价族 \cite{wu2022qpep} 为
\begin{equation}
\begin{split}
J_\theta(q,t)={}&v_2(q)^\top M_0(\theta)v_2(q)+2t^\top D(\theta)v_2(q)\\
&+t^\top G(\theta)t,
\end{split}
\label{eq:cost}
\end{equation}
满足 $q^\top q=1$ 且 $G(\theta)\succ0$。消去 $t$ 得到 $t^\star(q)=-G^{-1}Dv_2(q)$ 以及
\begin{equation}
\begin{gathered}
J_{\mathrm{red}}(q;\theta)=v_2(q)^\top M(\theta)v_2(q),\\
M=M_0-D^\top G^{-1}D .
\end{gathered}
\label{eq:reduced}
\end{equation}
映射 $\theta\mapsto(M_0,D,G)$ 假设为 $C^2$（定理~\ref{thm:implicit}）或 $C^3$（定理~\ref{thm:laplace}）。优化变量是四维的；数据只通过固定尺寸的 $10\times10$ 矩阵 $M$ 进入，该矩阵是每次观测组装一次的加权 Gram 矩阵。由于 $v_2$ 是 $\mathrm{Sym}^2(\mathbb{R}^4)$ 的一组基，任何关于 $R(q)$ 各元素为二次的四次型都可以写成 \eqref{eq:reduced} 的形式；这就是该理论所覆盖的函数类。

在驻点 $q^\star$ 处有 $\nabla J=\lambda q$，并且由欧拉齐次性有 $\lambda=4J$（定理~\ref{thm:implicit}）。全文使用的切基为
\begin{equation}
\begin{gathered}
E(q)=\begin{bmatrix}-q_1&-q_2&-q_3\\ q_0&-q_3&q_2\\ q_3&q_0&-q_1\\ -q_2&q_1&q_0\end{bmatrix},\\
q^\top E=0,\qquad E^\top E=I_3,
\end{gathered}
\label{eq:tangent}
\end{equation}
于是切向扰动为 $\delta q=E\,\delta\xi$。对应的旋转向量为 $\omega=2\,\delta\xi$；因此旋转向量坐标下的 Hessian 为 $H_\omega=\tfrac14H_{\mathbb{S}^3}$，即有一个因子 $4$ 把我们的曲率与按单位旋转角表示的曲率分开。QPEP 的可观测性与协方差分析是在其自身参数化下定义的，比较时应按此常数换算。

\textbf{假设。}
\begin{enumerate}
\item[(A1)] 对所考虑的论域中所有 $\theta$，$G(\theta)\succ0$ 且 $M(\theta)\in\mathbb{S}^{10}$。这使 $t^\star$ 唯一，并使约化代价成为球面上定义良好的四次型。
\item[(A2)] 前向代数求解器返回 $\nabla J=4Jq$、$q^\top q=1$ 的\emph{全部}实 KKT 根，直至符号对 $\pm q$。A2 是前向求解器的性质，而不是理论的性质；第~\ref{sec:rootcount} 节测量该系统的通用复根计数，第~\ref{sec:risks} 节说明实验中所用枚举器尚未被验证的部分。
\end{enumerate}

该层或者返回确定性的全局位姿 $(q^{(1)},t^\star(q^{(1)}))$，或者返回枚举极小值 $\{(q^{(k)},w_k)\}_{k=1}^{K}$ 上的平滑后验以及相应的位姿估计。我们记 $\phi=[\mathrm{vec}\,R(q);t^\star(q)]\in\mathbb{R}^{12}$ 为导数经过数值校验的位姿向量，$\Delta J=J(q^{(2)})-J(q^{(1)})$ 为最佳与次佳不同 $SO(3)$ 类之间的目标间隙，$K$ 为 $J_{\mathrm{red}}$ 在 $\mathbb{S}^3/\{\pm1\}$ 上被枚举出的真实局部极小值个数。本文所有旋转误差都是真旋转上的测地角 $\arccos((\operatorname{tr}(R_{\rm est}^\top R_{\rm gt})-1)/2)$，单位为度；所有平移误差都是欧氏误差，单位为米。
